\documentclass{osucourse}
\course{4350}

\begin{document}

\CatalogDescription
\PrerequisitesAndExclusions

\begin{Purpose}
Mathematical models and computational methods have been very useful in
understanding biological mechanisms underlying neuronal behavior. The
Hodgkin‐Huxley model, for example, has formed the basis for our
understanding of how action potentials are generated and how they
propagate along a nerve axon. More recently, mathematical models have
been used to help understand cellular processes responsible for both
normal and pathological firing patterns that arise in a wide range of
neuronal systems. Examples include models for sensory processing, motor
control, neurological disease, sleep rhythms and working memory.

This course provides a detailed introduction to how mathematical and
computational methods have been used to both develop and analyze models
that arise in neuroscience. We begin by deriving the Hodgkin‐Huxley
model and then describe dynamical system methods for analyzing models.
After discussing the dynamics of single neurons, we consider neuronal
networks and describe how different types of population firing patterns
depend on biological details, such as the intrinsic properties of
individual neurons and synaptic coupling. We conclude by considering
specific brain systems.
\end{Purpose}

\begin{Text}
\emph{Foundations of Mathematical Neuroscience}, by G. Bard Ermentrout
and David H. Terman
\end{Text}

\begin{TopicsList}
\item
  Overview: Neurons, synapses, neuronal firing patterns
\item
  Hodgkin‐Huxley Model: Resting potential, Nernst equation,
  Goldman‐Hodgkin‐Katz equation, cable equation, action potential
\item
  Dynamics I: Introduction to differential equations; phase‐planes;
  oscillations
\item
  Dynamics II: Stability analysis, bifurcation theory, numerical methods
\item
  Single cell dynamics I: Propagating action potentials; rhythmic
  behavior
\item
  Single cell dynamics II: Variety of channels, bursting oscillations;
  dendrites; multi‐compartment models
\item
  Synapses: Simple networks
\item
  Networks: Classification of network behavior; synchrony, role of
  different types of channels and coupling
\item
  Specific brain systems: Possible topics include models for working
  memory, vision, olfaction, sleep, Parkinson's tremor, stroke
\item
  Presentation of projects
\end{TopicsList}

\begin{SuggestedSchedule}
\week{1}{Overview: neurons, synapses, firing patterns.}
\week{2}{Resting potential; Nernst and Goldman-Hodgkin-Katz equations.}
\week{3}{Cable equation; derivation of the Hodgkin-Huxley model.}
\week{4}{Differential equations, phase planes, oscillations.}
\week{5}{Stability analysis and bifurcation theory.}
\week{6}{Numerical methods; simulation of model neurons.}
\week{7}{Propagating action potentials.}
\week{8}{Rhythmic behavior; midterm.}
\week{9}{Channel varieties; bursting oscillations.}
\week{10}{Dendrites and multi-compartment models.}
\week{11}{Synapses; simple networks.}
\week{12}{Classification of network behavior; synchrony and coupling.}
\week{13}{Specific brain systems: working memory, vision, olfaction.}
\week{14}{Specific brain systems: sleep, Parkinson's tremor, stroke.}
\finalsweek{Presentation of projects.}
\end{SuggestedSchedule}

\end{document}
